\scanpage{20}
\mspage{16}
% archive sheet number: 19

disjointe autoduale de
\[
\widehat\delta\Phi\overset{\mathrm{def}}=\delta\Phi\cup\partial I=\delta\Phi\cup\underbrace{(\partial I\setminus\unsure{\delta I|\partial I}\cap\delta\Phi)}_{=\Phi^*_{\text{sév}}}
\]
sév = semi-évanescents, ainsi.
\diagram{\begin{tikzcd}[ampersand replacement=\&,column sep=small,row sep=large]
\widehat\delta\Phi=\Phi_{\mathrm{is}} \arrow[rr,leftrightarrow,bend right=40] \arrow[rr,leftrightarrow,bend right=44] \arrow[rr,bend right=48] \& {}\cup{} \& \Phi_{\mathrm{lac}} \& {}\cup\Phi_{\mathrm{ev}} \arrow[d] \& {}\cup\Phi_{\mathrm{sev}} \arrow[r,leftrightarrow,bend right=30] \arrow[d,equal] \& {}\cup\Phi^{\unsure{*|'}}_{\mathrm{sev}} \arrow[d,equal] \& {}\cup\partial^*\Phi \arrow[loop below,leftrightarrow]\!\\
\&\&\&\substack{\text{disparaît}\\\text{par dualité}}\&\deleted{\unsure{\partial I\setminus\partial I\cap\delta\Phi^*}}\arrow[d,equal]\&\partial I\setminus\partial I\cap\delta\Phi\&\\
\&\&\&\&\deleted{\unsure{\delta\Phi\setminus\delta\Phi^*\cap\delta\Phi}}\&\&
\end{tikzcd}}
où
\[
\begin{aligned}
\Phi_{\mathrm{is}}&=\Phi^*_{\mathrm{lac}}\\
\Phi_{\mathrm{lac}}&=\Phi^*_{\mathrm{is}}\\
\partial^*\Phi_{\deleted{\unsure{\mathrm{sev}}}}&=\partial^*\Phi^*\\[1ex]
\Phi^*_{\mathrm{ev}}&=\Phi^{**}_{\mathrm{ev}}=\varnothing
\end{aligned}
\]

Pour les $s\in\partial I$, il y a quatre positions possibles p.r. à $\Phi$:
\[
\begin{aligned}
\text{a)}\quad s\in\Phi_{\mathrm{is}}&\Longleftrightarrow s\in\Phi^*_{\mathrm{lac}}\\
\text{b)}\quad s\in\Phi_{\mathrm{sev}}&\Longleftrightarrow s\notin\delta\Phi^*\\
\text{a$^*$)}\quad s\in\Phi_{\mathrm{lac}}&\Longleftrightarrow s\in\Phi^*_{\mathrm{is}}\\
\text{b$^*$)}\quad s\notin\delta\Phi&\Longleftrightarrow s\in\Phi^*_{\mathrm{sev}}
\end{aligned}
\]
\[
\begin{aligned}
\Phi_{\text{sév}}&=\{s\in\partial I\mid s\notin\delta\Phi^*\}.\\
&=\partial I\setminus\partial I\cap\delta\Phi^*\\
\Phi^*_{\mathrm{sev}}&=\{s\in\partial I\mid s\notin\delta\Phi\}\\
&=\partial I-\partial I\cap\delta\Phi
\end{aligned}
\]
% slanted
\margin{\uline{Cor 1} \deleted{\unsure{Omb}} $\unsure{\operatorname{Supp}^{\circ}|\operatorname{S-m}^{\circ}}(\Phi)=\operatorname{Supp}^{\circ}(\operatorname{Omb}^{\circ}\Phi)=$\par
\uline{Cor. 2} $\operatorname{Supp}^{\circ}\Phi=\operatorname{Omb}^{\circ}\Phi^{**}$}

On trouve alors
\[
\boxed{\unsure{\operatorname{cosupp}^{\circ}|\operatorname{cos-m}^{\circ}}(\Phi)=\operatorname{cosupp}^{\circ}(\operatorname{Omb}^{\circ}(\Phi))=\operatorname{Omb}^{\circ}(\Phi^*)}
\]

On peut appeler \uline{parties constructibles de $\unsure{M|\mathcal M}$} les parties de la forme $\operatorname{Omb}^{\circ}(\Phi)$, $\Phi$ polygone. \deleted{Il y a} L'application
\[
\Phi\longmapsto\operatorname{Omb}^{\circ}(\Phi)\qquad\operatorname{Polyg}(M)\longrightarrow\operatorname{Cons}(M)
\]
est bijective. En effet, on récupère $\Phi$
